\SourceLeaf{191}
\[
Pm=A+B\mathord{.}\varphi m
\]
\noindent wo
\[
48A=19\rho^{\frac{1}{2}}+7\rho,\qquad48B=5\rho^{-\frac{1}{2}}+17+19\rho^{\frac{1}{2}}+7\rho
\]
\noindent und nahe genug
\[
A=\tfrac{13}{24}\rho^{\frac{33}{52}},\qquad B=\rho^{\frac{7}{24}}
\]
\noindent für \(3\tfrac{1}{2}\) Proc.\quad \(\log\tfrac{1}{\rho}=0.0149403\)\quad Genau
\[
A=0.52998474,\qquad\log B=9.9956901
\]
\noindent die Näherungsformel gibt
\[
A=0.5299694,\qquad\log B=9.9956424
\]
\noindent Gerechnet war nach \(\tfrac{13}{24}+\varphi m=P^0\)\par
\noindent man kann also setzen
\[
\begin{aligned}
P&=\rho^{\frac{7}{24}}P^0-\tfrac{13}{24}(\rho^{\frac{7}{24}}-\rho^{\frac{3}{4}})\\
&=P^0-\tfrac{7}{24}(1-\rho)P^0-\tfrac{143}{576}(1-\rho)
\end{aligned}
\]
\noindent also für \(3\tfrac{1}{2}\) Proc.\quad \(P=P^0-\tfrac{49}{4968}P^0-\tfrac{1091}{119232}\)\par
\noindent für 4 Proc.\quad \(P=P^0-\tfrac{7}{624}P^0-\tfrac{11}{1152}\)

\Needspace{14\baselineskip}
Beispiel: Nr. 65. H.\ . 4 Proc.\quad \(m=53.784\)
\begin{center}
\fontsize{8}{13}\selectfont\setlength{\tabcolsep}{2pt}
\begin{tabular}{ll}
\(\log\varphi(m-\tfrac{13}{24})\ .\ .\ .\ .\ .\)&1.05488\\
\(\text{comp.}\log\rho^{\frac{1}{4}}\ .\ .\ .\ .\ .\ .\)&0.00426\\
\(\log g(m+\tfrac{11}{49})^{\frac{13}{28}}\ .\ .\ .\ .\)&0.00440\\\cline{2-2}
&1.06354\\
\(P=\)&11.5755
\end{tabular}\qquad
\begin{tabular}{ll}
\(\log\varphi m\ .\ .\ .\ .\)&1.04830\\
\(\log\rho^{\frac{7}{24}}\ .\ .\)&\(-0.00497\)\\
\(B\varphi m\ .\ .\ .\ .\)&11.0492\\
\(\tfrac{13}{24}\rho^{\frac{3}{4}}\ .\ .\ .\ .\ .\)&0.5260\\\cline{2-2}
&11.5752
\end{tabular}\par\medskip
Verbesserung des frühern Werthes \(P^0\)\par\smallskip
\(P^0-\tfrac{1}{99}\tfrac{105}{104}P^0-\tfrac{11}{1152}\)\par\smallskip
\begin{tabular}{r}
\(\tfrac{13}{24}+\varphi m=11.7179\)\\
\(-0.1314\ 5\)\\
\(-0.0095\ 5\)\\\hline
11.5769
\end{tabular}
\end{center}

\Needspace{5\baselineskip}
\pdfbookmark[2]{Stehende Ehe.}{witwen-stehende-ehe}
\begin{center}\textit{Stehende Ehe.}\end{center}
\noindent\hspace*{1.2em} Alter des Mannes \(M\), der Frau \(m\) zur Zeit von Oct. 1. Sterblichkeitstafel lebende Männer vom Alter \(x=Fx\), Werth der Witwenpension \(=R-Q\).

[Nach den Statuten nimmt die Pension mit dem Ablaufe des Gnadenquartals ihren Anfang, wird in halbjährigen am 1. April und am 1. October jeden Jahres fälligen Raten ausgezahlt und erlischt mit dem Schlusse des Sterbemonats. Für den Zeitwerth \(R-Q\) einer etwa eintretenden Witwenpension welche jährlich die Münzeinheit beträgt ist demnach
\[
\begin{aligned}
(R-Q)\mathord{.}fm\mathord{.}Fm&=\tfrac{1}{12}\rho\{Fm-F(m+\tfrac{1}{4})\}\\
&\qquad\{f(m+\tfrac{6}{12})+f(m+\tfrac{7}{12})+f(m+\tfrac{8}{12})\}\\
&\quad+\tfrac{1}{12}\rho\{Fm-F(m+\tfrac{2}{4})\}\\
&\qquad\{f(m+\tfrac{9}{12})+f(m+\tfrac{10}{12})+f(m+\tfrac{11}{12})\}\\
&\quad+\tfrac{1}{12}\rho^{\frac{3}{2}}\{Fm-F(m+\tfrac{3}{4})\}\\
&\qquad\{f(m+\tfrac{12}{12})+f(m+\tfrac{13}{12})+f(m+\tfrac{14}{12})\}\\
&\quad+\tfrac{1}{12}\rho^{\frac{3}{2}}\{Fm-F(m+\tfrac{4}{4})\}\\
&\qquad\{f(m+\tfrac{15}{12})+f(m+\tfrac{16}{12})+f(m+\tfrac{17}{12})\}\\
&\quad+\text{u.\,s.\,w.}
\end{aligned}
\]
\noindent oder:]
\[
\begin{aligned}
R\mathord{.}fm&=\tfrac{1}{2}\rho f(m+\tfrac{17}{24})+\tfrac{1}{2}\rho^{\frac{3}{2}}f(m+\tfrac{29}{24})+\tfrac{1}{2}\rho\rho f(m+\tfrac{41}{24})+.\ .\\
&=\rho^{\frac{5}{4}}f(m+\tfrac{23}{24})+\rho^{\frac{9}{4}}f(m+\tfrac{47}{24})+.\ .
\end{aligned}
\]
\[
R=\tfrac{1}{24}\rho^{\frac{5}{4}}+\rho^{\frac{7}{24}}\varphi m
\]
\[
\begin{aligned}
Q\mathord{.}fm\mathord{.}Fm&=\rho^{\frac{5}{4}}f(m+\tfrac{23}{24})\mathord{.}F(M+\tfrac{5}{8})\\
&\quad+\rho^{\frac{9}{4}}f(m+\tfrac{47}{24})F(M+\tfrac{13}{8})+.\ .
\end{aligned}
\]
